% ===== §7 The Aesthetic Perception of Contingency =====
\section{The Aesthetic Perception of Contingency}
\label{sec:aesthetic-perception}

\noindent
The construction of the field is an aesthetic problem; so, it turns out, is the perception of what arises within it. This section turns from the making of the field to the catching of what the field lets in, and argues that the capacity to catch it is itself an aesthetic capacity, the second of the paper's named concepts: the \emph{aesthetic perception of contingency}. The argument also lightens, in an important way, the apparent burden of the previous section. For if constructing fields were the only path to the good this paper describes, that good would be effortful and rare, available chiefly to those with the leisure and means to build elaborate frames. This section shows that the deeper capacity is lighter and more portable than the construction of fields: it is a perceptual skill, exercisable anywhere, that catches the contingent beauty the world is always offering and turns it toward the beloved. A wildflower at the roadside is enough.

% ============================================================
\subsection{Contingency is everywhere; the problem is catching it}
\label{sub:ap-everywhere}

The previous sections may have given the impression that contingency must be manufactured, that the everyday is a desert of pure repetition into which openness has to be deliberately imported by the construction of fields. This impression is false, and correcting it is the starting point of the present section. The everyday is not, in fact, pure repetition; its repetitions are shot through with contingency at every moment, with small undetermined irruptions that the routine does not fully suppress. A wildflower has come up at the edge of the path. The light has fallen, just now, in a particular way across a wall. A stranger's face, in passing, has worn an expression that means something. A phrase overheard, a sound, a sudden smell, the everyday leaks contingency continuously, through the seams of its repetition, in a steady stream of small undetermined events. The problem is not that the everyday lacks contingency. The problem is that almost all of this contingency goes \emph{uncaught}, arrives and departs unnoticed, unregistered, unturned-toward, because the purposive perception of the everyday, fixed on its functions and its destinations, is blind to it.

This reframes the whole question of practice. The deepest practical capacity is not, after all, the construction of fields, demanding as that is; it is the capacity to \emph{catch} the contingency that is already arriving, to notice and receive what the everyday is already offering through its seams. And the minimal event in which this capacity is exercised is very small and entirely free: one catches the wildflower at the roadside, sees it, is touched by it, and turns it toward the one beside one: \emph{look}. In that small act, a contingency that would have passed unnoticed is caught, and in being turned-toward it is shared, and in being shared it enters the structure of the ``we'' and becomes a deposit, however tiny, in the relation's wealth. The flower, caught and shared, is a contingency that has intervened in the relational structure, the smallest possible instance of everything this paper has described, requiring no constructed field, no journey, no leisure, nothing but the capacity to catch what is already there and turn it toward the beloved. The chain runs: contingency \(\to\) perception \(\to\) sharing \(\to\) wealth, and it can run anywhere, at any moment, for free.

% ============================================================
\subsection{The three moments of the perception}
\label{sub:ap-moments}

The capacity so described is not simple; it has three distinct moments, each a distinct ability that can be present or absent, developed or atrophied, and it is worth separating them because the cultivation of the capacity (\xref{sec:aesthetic-education}) is the cultivation of each.

The first moment is \emph{recognition}: actually \emph{seeing} the contingent thing when it arrives. This is harder than it sounds, and most contingency fails at exactly this first step. The purposive perception of the everyday, the perception that sees the road but not the flower, because it is fixed on getting somewhere, is structurally blind to the contingent, since the contingent is by definition what is not on the way to the purpose. To see the flower is to perceive in a different mode: a mode in which purpose is suspended and the perception is open to what is there for its own sake rather than for its use. This is the mode classically called \emph{aesthetic perception}, the disinterested attention that beholds a thing for what it is rather than for what it is for, that sees the flower as a flower and not as an irrelevance beside the path to the destination. The first moment of catching contingency is thus an aesthetic suspension of the purposive gaze, the opening of perception to the gratuitous, the un-useful, the merely there.

The second moment is \emph{being moved}: not merely seeing the contingent thing but being \emph{touched} by it, letting it stir something. Seeing is not enough; one can register a flower coldly, as a fact, and be unmoved. The second moment is the receptivity in which the seen thing is allowed to matter, to provoke a small movement of feeling, to be received as significant rather than merely noted. This is a poetic receptivity, and the Chinese tradition has a precise name for it: \emph{xing} (\zh{兴}), the stirring of feeling by a thing, the rising of emotion or significance in response to a perceived object, which the classical poetics placed at the very root of poetry, ``the poem begins in \emph{xing},'' in the affective stirring that a thing provokes \citep{owen1992}. To be moved by the roadside flower is to undergo \emph{xing}: to let the small contingent thing stir something, to be affectively reached by it rather than merely registering its presence. This receptivity, too, can be developed or dulled; a life of purely purposive perception dulls it, and the practiced sensibility of the poet keeps it alive.

The third moment is \emph{turning-to-share}: taking what one has seen and been moved by, and turning it toward the other. \emph{Look.} This small act is the one that makes the whole chain relational, and it is the decisive moment for the purposes of this paper, because it is what converts a private aesthetic experience into a relational event. One could see the flower and be moved by it and keep it, a complete solitary aesthetic experience, of the kind the lone traveller has. The turning-toward is what makes it a deposit in the ``we'': by sharing the caught contingency, one lets it intervene in the relational structure, makes it a moment of the ``we'' rather than of the ``I.'' The turning-to-share is the relational moment of the aesthetic perception of contingency, and it is what distinguishes the capacity this paper describes from the solitary aesthetic sensibility the tradition has long praised. The full capacity is the chain of all three: to see the contingent (recognition), to be moved by it (\emph{xing}), and to turn it toward the beloved (sharing), and these three, as the philosophy of beauty will show (\xref{sub:ap-antinomy}), are the three registers brought into accord: the real kernel of the contingent thing encountered in recognition, the imaginary investment of being-moved, and the symbolic articulation of the turning-toward. On that account the \emph{look} is not the aftermath of a completed aesthetic judgement but its very completion, the moment at which the relational subject's judgement achieves the shared symbolic accord its relational constitution demands, and a relationship rich in this capacity, in which both partners habitually catch and share the contingent beauty the world leaks, is rich in a continuous stream of small generative deposits that require no journey and no constructed field, only the cultivated perception that catches what is already there.

% ============================================================
\subsection{A skill, not an engineering: the portable container}
\label{sub:ap-skill}

This resolves, at last, the tension that the construction of fields had introduced, the worry that the good of this paper is effortful, elaborate, and available only to those who can build grand frames. The aesthetic perception of contingency is not an engineering but a \emph{skill}: not a structure one builds in the world but a capacity one carries in oneself, exercisable anywhere, requiring no external apparatus. And a skill, unlike a constructed field, is \emph{portable}. The constructed field is a scaffolding, an external structure, built at some cost in time and means, that holds open a space for the relation to come to presence. The aesthetic perception of contingency is the internalisation of what the scaffolding was for: once one has the skill, one no longer needs the scaffolding, because one carries within oneself the capacity that the scaffolding was constructed to supply. The skilled perceiver opens the field wherever they are, catches the contingent beauty the everyday leaks, turns it toward the beloved, and so generates, in the midst of the most ordinary day, the small relational events that the constructed field was built to produce.

This is the movement from \emph{scaffolding to skill}, and it is the deep logic of the paper's eventual generalisation of travel (\xref{sec:generalized}). The literal journey is a scaffolding: a powerful, expensive, external frame that opens the field all at once, and whose chief value, beyond the field it opens, is that it \emph{teaches}, it lets one experience the field vividly enough to recognise it, to learn what it is and what it feels like to catch and share the contingent within it. But the mature form of the practice is not the perpetual construction of scaffoldings, which would indeed be effortful and elite; it is the skill that the scaffolding taught, carried back into the everyday, exercisable without the scaffolding. The journey shows what the field is; the skill lets one open it anywhere. And spaciousness is what makes the catching possible at every point: the skilled perceiver is one who keeps, within themselves and within their day, enough emptiness, enough freedom from the total occupation of purpose, that the contingent has room to be noticed when it arrives. Without that interior spaciousness the contingency leaks past uncaught, the perception too full of purpose to register the flower. The skill, then, is inseparable from the cultivation of an interior emptiness: the aesthetic perception of contingency is, among other things, the capacity to keep oneself empty enough to catch what comes.

The skill connects, finally, to two things the series has already established. It is the positive, capacitated form of what the previous paper marked as the un-programmability of the good \citep{paperxx}: there it was shown that the good cannot be reduced to a procedure and must be judged afresh in each case; here that un-programmable judgement is given its positive name as a cultivated perceptual skill, the thing one develops where a procedure cannot be supplied. And it is the everyday exercise of what the series called witnessing \citep{paperxix}: to catch the flower and say \emph{look} is the smallest act of shared witnessing, the turning of two gazes toward one thing, the constitution of a shared seeing that is a moment of the ``we.'' The aesthetic perception of contingency is witnessing made into a daily skill, and its cultivation, to which the paper turns next, is the cultivation of the capacity to witness the world together, beautifully, at the scale of a roadside flower.

% ============================================================
\subsection{Eastern resources: the poetics of the wayfarer and the landscape}
\label{sub:ap-eastern}

The capacity this section describes is not without its masters, and the richest tradition of them is the East Asian poetics of the wayfarer and the landscape, which developed the aesthetic perception of contingency to a height the Western literature of travel rarely reached, and which the paper draws on both for its phenomenological precision and as a counterweight to the otherwise Western cast of its aesthetics. Three strands are especially apposite.

The first is the travel poetics of Bash\=o, whose \emph{Oku no Hosomichi}, the narrow road to the deep north, is the supreme record of the aesthetic perception of contingency exercised on a journey \citep{basho1966}. Bash\=o's practice is exactly the chain this section has described: the catching of small contingent things (a particular silence, the cry of a cicada, the look of a ruined gate), the being-moved by them (the \emph{haiku} is the deposit of a \emph{xing}, the stirring crystallised), and in the form of the travel diary itself, the turning of them toward another, the sharing of the caught moment with a reader and, within the journey, with the companion Sora who travelled beside him. Bash\=o's aesthetic of \emph{sabi} and the lightness he called \emph{karumi} are, among other things, disciplines of the perception of contingent beauty, trained capacities to catch and be moved by the small, the passing, the un-grand, exactly the contingencies the purposive gaze misses.

The second is the Chinese tradition of the landscape travel-record and the landscape poem, the \emph{youji} of writers from Liu Zongyuan to Xu Xiake, and the landscape poetry of Xie Lingyun, Wang Wei, and their lineage, in which the journey into mountains and waters is at once a literal travel and a discipline of perception, the cultivation of an eye and a heart attuned to the beauty of the contingent and the spacious \citep{owen1996}. The Chinese landscape aesthetic is, in a deep sense, an aesthetic of spaciousness: the valuing of emptiness in the painting and the poem, the significance of what is not depicted, the use of the void, the unbrushed silk, the unsaid word, the empty distance, as the very medium of beauty. This is the aesthetic correlate of the eleventh chapter's metaphysics of emptiness (\xref{sub:field-spaciousness}): a whole tradition of perception trained to find the use of the vessel in its emptiness, the beauty of the landscape in its spacious void.

The third is the broader figure, common to the tradition, of travel as a discipline of perception and not merely a movement through space, the wayfarer whose journey is an education of the eye and the heart, who travels in order to learn to perceive. This figure gives the paper's eventual generalisation an ancient precedent: the point of the journey was never the destinations but the trained perception the journey cultivated, the eye for contingent and spacious beauty that, once trained, could be exercised anywhere. The masters of this tradition are, in the paper's terms, those who completed the movement from scaffolding to skill, who used the journey to cultivate a perception that they then carried into all of life, so that for them the whole world became a landscape to be perceived with the wayfarer's trained and open eye. They are the proof that the capacity this section describes can be cultivated to mastery, and the bridge to the question of how such cultivation is done, the question of aesthetic education, to which the paper now turns.
